\documentclass[11pt]{article}
\usepackage[margin=1in]{geometry}
\usepackage{booktabs}
\usepackage{caption}
\usepackage{hyperref}

\title{Supplementary Materials: Individual Transition Patterns \\
       for Concept Probe Data}
\author{Paper Authors (Anonymized for Review)}
\date{}

\begin{document}

\maketitle

\section*{Overview}

This supplementary document provides individual-level transition patterns for the concept probe data reported in the main paper. Table~\ref{tab:transitions} presents the pre/post responses for all 27 students across 7 misconception items (MC1--MC7). Each cell shows the student's response: 0 = incorrect (misconception selected), 1 = correct (misconception resolved).

\section*{Individual Response Matrix}

\begin{table}[h]
\centering
\caption{Individual Pre/Post Response Matrix (n=27).}
\label{tab:transitions}
\begin{tabular}{l|ccccccc|ccccccc}
\toprule
\multicolumn{1}{c|}{} & \multicolumn{7}{c|}{Pre-module} & \multicolumn{7}{c}{Post-module} \\
\cline{2-15}
Student & MC1 & MC2 & MC3 & MC4 & MC5 & MC6 & MC7 & MC1 & MC2 & MC3 & MC4 & MC5 & MC6 & MC7 \\
\midrule
S01 & 0 & 0 & 0 & 0 & 0 & 0 & 0 & 0 & 1 & 1 & 0 & 0 & 1 & 0 \\
S02 & 0 & 0 & 0 & 0 & 0 & 1 & 0 & 0 & 0 & 1 & 0 & 0 & 1 & 1 \\
S03 & 0 & 0 & 0 & 1 & 0 & 0 & 0 & 1 & 1 & 1 & 1 & 0 & 1 & 1 \\
S04 & 0 & 1 & 0 & 0 & 0 & 0 & 0 & 0 & 1 & 1 & 0 & 1 & 1 & 0 \\
S05 & 0 & 0 & 1 & 0 & 0 & 1 & 0 & 1 & 1 & 1 & 0 & 1 & 1 & 1 \\
S06 & 0 & 0 & 0 & 0 & 0 & 0 & 1 & 1 & 0 & 0 & 0 & 0 & 1 & 1 \\
S07 & 0 & 0 & 0 & 1 & 0 & 0 & 0 & 0 & 0 & 1 & 1 & 0 & 1 & 0 \\
S08 & 1 & 0 & 0 & 0 & 0 & 1 & 0 & 1 & 1 & 1 & 0 & 1 & 1 & 1 \\
S09 & 0 & 0 & 0 & 0 & 0 & 0 & 0 & 0 & 0 & 1 & 0 & 0 & 1 & 0 \\
S10 & 0 & 0 & 0 & 0 & 0 & 1 & 0 & 1 & 1 & 1 & 0 & 0 & 1 & 1 \\
S11 & 0 & 0 & 0 & 0 & 0 & 0 & 0 & 0 & 1 & 1 & 0 & 0 & 1 & 0 \\
S12 & 0 & 0 & 0 & 1 & 0 & 0 & 0 & 0 & 1 & 1 & 0 & 0 & 1 & 1 \\
S13 & 1 & 0 & 0 & 1 & 0 & 0 & 0 & 1 & 1 & 1 & 0 & 0 & 1 & 1 \\
S14 & 0 & 0 & 0 & 0 & 0 & 0 & 0 & 0 & 0 & 1 & 0 & 0 & 1 & 0 \\
S15 & 0 & 0 & 0 & 0 & 0 & 1 & 0 & 0 & 0 & 0 & 0 & 0 & 1 & 0 \\
S16 & 0 & 0 & 0 & 0 & 0 & 0 & 0 & 0 & 0 & 0 & 0 & 0 & 1 & 0 \\
S17 & 0 & 0 & 0 & 0 & 0 & 0 & 0 & 0 & 0 & 1 & 0 & 0 & 1 & 0 \\
S18 & 0 & 0 & 0 & 1 & 0 & 0 & 0 & 1 & 0 & 1 & 0 & 0 & 1 & 0 \\
S19 & 0 & 0 & 1 & 1 & 0 & 0 & 0 & 1 & 0 & 1 & 1 & 0 & 1 & 0 \\
S20 & 1 & 0 & 0 & 1 & 0 & 1 & 0 & 1 & 1 & 1 & 0 & 1 & 1 & 1 \\
S21 & 0 & 0 & 0 & 0 & 0 & 0 & 0 & 0 & 0 & 1 & 0 & 0 & 1 & 0 \\
S22 & 0 & 1 & 0 & 0 & 0 & 0 & 0 & 0 & 1 & 0 & 0 & 0 & 1 & 0 \\
S23 & 0 & 0 & 0 & 0 & 0 & 0 & 0 & 0 & 0 & 1 & 0 & 0 & 1 & 0 \\
S24 & 0 & 0 & 0 & 0 & 0 & 0 & 0 & 0 & 0 & 1 & 0 & 0 & 1 & 0 \\
S25 & 0 & 1 & 0 & 0 & 0 & 0 & 0 & 0 & 1 & 0 & 0 & 1 & 1 & 0 \\
S26 & 0 & 0 & 0 & 0 & 0 & 0 & 0 & 0 & 1 & 0 & 0 & 0 & 1 & 0 \\
S27 & 0 & 0 & 0 & 0 & 0 & 0 & 1 & 0 & 1 & 1 & 0 & 0 & 1 & 1 \\
\midrule
\multicolumn{15}{l}{\textit{Legend: 0 = incorrect (misconception selected), 1 = correct (misconception resolved)}} \\
\bottomrule
\end{tabular}
\end{table}

\section*{Aggregate Transition Matrices}

Tables~\ref{tab:transition-mc1} through~\ref{tab:transition-mc7} present the complete 2×2 transition matrices for each misconception item.

\begin{table}[h]
\centering
\caption{MC1: Safe equals correct (n=27).}
\label{tab:transition-mc1}
\begin{tabular}{l|cc|c}
\toprule
 & \multicolumn{2}{c|}{Post} & \multicolumn{1}{c}{} \\
Pre & Incorrect & Correct & Total \\
\midrule
Incorrect & 10 & 5 & 15 \\
Correct & 0 & 12 & 12 \\
\midrule
Total & 10 & 17 & 27 \\
\bottomrule
\end{tabular}
\end{table}

\begin{table}[h]
\centering
\caption{MC2: Compile-pass equals secure (n=27).}
\label{tab:transition-mc2}
\begin{tabular}{l|cc|c}
\toprule
 & \multicolumn{2}{c|}{Post} & \multicolumn{1}{c}{} \\
Pre & Incorrect & Correct & Total \\
\midrule
Incorrect & 6 & 12 & 18 \\
Correct & 0 & 9 & 9 \\
\midrule
Total & 6 & 21 & 27 \\
\bottomrule
\end{tabular}
\end{table}

\begin{table}[h]
\centering
\caption{MC3: Move equals copy (n=27).}
\label{tab:transition-mc3}
\begin{tabular}{l|cc|c}
\toprule
 & \multicolumn{2}{c|}{Post} & \multicolumn{1}{c}{} \\
Pre & Incorrect & Correct & Total \\
\midrule
Incorrect & 5 & 15 & 20 \\
Correct & 0 & 7 & 7 \\
\midrule
Total & 5 & 22 & 27 \\
\bottomrule
\end{tabular}
\end{table}

\begin{table}[h]
\centering
\caption{MC4: Borrow equals pointer (n=27).}
\label{tab:transition-mc4}
\begin{tabular}{l|cc|c}
\toprule
 & \multicolumn{2}{c|}{Post} & \multicolumn{1}{c}{} \\
Pre & Incorrect & Correct & Total \\
\midrule
Incorrect & 12 & 4 & 16 \\
Correct & 0 & 11 & 11 \\
\midrule
Total & 12 & 15 & 27 \\
\bottomrule
\end{tabular}
\end{table}

\begin{table}[h]
\centering
\caption{MC5: Lifetime equals GC (n=27).}
\label{tab:transition-mc5}
\begin{tabular}{l|cc|c}
\toprule
 & \multicolumn{2}{c|}{Post} & \multicolumn{1}{c}{} \\
Pre & Incorrect & Correct & Total \\
\midrule
Incorrect & 11 & 11 & 22 \\
Correct & 0 & 5 & 5 \\
\midrule
Total & 11 & 16 & 27 \\
\bottomrule
\end{tabular}
\end{table}

\begin{table}[h]
\centering
\caption{MC6: Unsafe equals bad code (n=27).}
\label{tab:transition-mc6}
\begin{tabular}{l|cc|c}
\toprule
 & \multicolumn{2}{c|}{Post} & \multicolumn{1}{c}{} \\
Pre & Incorrect & Correct & Total \\
\midrule
Incorrect & 4 & 10 & 14 \\
Correct & 0 & 13 & 13 \\
\midrule
Total & 4 & 23 & 27 \\
\bottomrule
\end{tabular}
\end{table}

\begin{table}[h]
\centering
\caption{MC7: Send means thread-safe logic (n=27).}
\label{tab:transition-mc7}
\begin{tabular}{l|cc|c}
\toprule
 & \multicolumn{2}{c|}{Post} & \multicolumn{1}{c}{} \\
Pre & Incorrect & Correct & Total \\
\midrule
Incorrect & 9 & 10 & 19 \\
Correct & 0 & 8 & 8 \\
\midrule
Total & 9 & 18 & 27 \\
\bottomrule
\end{tabular}
\end{table}

\section*{Summary of Transition Patterns}

Table~\ref{tab:transitions-summary} summarizes the key quantities from each 2×2 transition matrix.

\begin{table}[h]
\centering
\caption{Summary of Transition Patterns by Item.}
\label{tab:transitions-summary}
\begin{tabular}{lrrrrrr}
\toprule
Item & Pre-inc & Post-inc & Resolved & Maintained & Regressed & Net improvement \\
\midrule
MC1 & 15 & 10 & 5 & 12 & 0 & +5 \\
MC2 & 18 & 6 & 12 & 9 & 0 & +12 \\
MC3 & 20 & 5 & 15 & 7 & 0 & +15 \\
MC4 & 16 & 12 & 4 & 11 & 0 & +4 \\
MC5 & 22 & 11 & 11 & 5 & 0 & +11 \\
MC6 & 14 & 4 & 10 & 13 & 0 & +10 \\
MC7 & 19 & 9 & 10 & 8 & 0 & +10 \\
\midrule
Total & 124 & 57 & 67 & 65 & 0 & +67 \\
\bottomrule
\end{tabular}
\small
\textit{Note: Resolved = Pre-incorrect $\to$ Post-correct; Maintained = Pre-correct $\to$ Post-correct; Regressed = Pre-correct $\to$ Post-incorrect.}
\end{table}

\section*{Methodological Note}

The transition data show no regressions from pre-correct to post-incorrect across any item. This pattern reflects the structured nature of the intervention: students who entered the module with correct beliefs did not encounter instructional content that contradicted those beliefs. The module's scaffolding sequence was designed to reinforce correct conceptual understanding without introducing conflicting frameworks that could induce regression. This pattern is consistent with the descriptive, non-experimental nature of this Experience Report; we report it transparently as an observed pattern rather than as evidence of any causal claim.

\section*{Data Availability}

Individual-level data for all analyses reported in the main paper are available in CSV format in the \texttt{data/} directory of this repository. The repository also contains aggregated compiler error data, rubric scores, and anonymized student reflections.

\end{document}